\chapter{Appendix} % (fold)
\label{appendix}

\begin{center}
  \textsc{Extracting chirps for model fitting in full simulations}
\end{center}

I used the following pipeline to extract estimates on how the instantaneous
frequency evolves during a chirp using the output of the sliding convolutional
classifier as a prior on the Raab dataset:

\begin{enumerate}

  \item Electrode selection: By the power of the frequency band of the EOD of
    the fish I extracted the signal on the electrode of the strongest power,
    i.e., the electrode that the fish was closest to. For each chirp, I took a
    full second, meaning 500 \si{\milli\second} before and after the chirp. The
    result is a single voltage trace over one second, i.e. 20,000 samples, as
    the signals were sampled at \SI{20}{\kilo\hertz}.

  \item Band-pass filtering: I then applied a band-pass filter after
    multiplying the snipped with a Tukey window ($\alpha=0.1$) to reduce edge
    effects of the filter. The cutoff frequency below the EOD frequency
    of the upper fish EOD$f_{\text{up}}$ was set relative to that of the lower fish EOD$f_{\text{low}}$ at $\text{EOD}f_{\text{low}} + 0.8 \times
    \Delta\text{EOD}f$ and above the EOD $f_{\text{up}}$ at $1.9
    \times$EOD$f_{\text{low}}$. The lower cutoff frequency was chosen to reduce
    interference from the EOD of the other fish while not cutting off potential
    down-modulations of the frequency that sometimes occur after chirps. The
    upper cutoff was chosen to reduce the interference from the first harmonic
    of the EOD of the lower fish while not cutting off the top of the chirp.

  \item Standardization: I then standardized the filtered EOD to have a mean of zero
    and a standard deviation of one.

  \item Instantaneous frequency: I then computed the instantaneous frequency of
    the signal using zero crossings. The instantaneous frequency was
    additionally smoothed using a Gaussian kernel with a standard deviation of
    5 samples. It was then linearly interpolated to the original sample rate of
    the signal. The result is an estimate on how the EOD$f_{up}(t)$ evolved
    over the course of one second EOD$f_{\text{inst}}(t)$ at the temporal
    resolution of a single period of the EOD of the upper fish.

  \item Centering baseline frequency: To center the baseline of this snippet of
    the instantaneous frequency that contains a chirp to zero, I estimated the
    baseline frequency from the largest peak of the distribution of all
    frequencies in the snippet. To estimate it, I used a histogram of 100 bins
    as this was much faster than using a kernel density estimate. The baseline
    frequency estimate was then subtracted from the instantaneous frequency to
    shift the baseline frequency to zero.

  \item Peak detection: The most descriminating feature of chirps in
    Apteronotids is the rapid increase of the EOD$f$ during a chirp. To
    estimate the exact time of the chirp, I detected peaks in the instantaneous
    frequency that exceeded the 99\th percentile of the centered instantaneous
    frequencies using the \texttt{find\_peaks} function of the \texttt{scipy}
    package \parencite{virtanenSciPyFundamentalAlgorithms2020}. Before
    detecting the peaks, I applied another Tukey window ($\alpha=0.3$) to the
    centered instantaneous frequencies to scale down fluctuations at the edges
    of the snippet (as chirps are expected in the center of the snippet). I
    only considered the largest peak and discarded the whole snippet to
    continue with the next chirp when the peak was more than 250
    \si{\milli\second} away from the center of the snippet. I also discarded
    the whole snippet if included multiple peaks that where larger than $0.3
    \times \max{\text{EOD}f_{\text{snip}}}$.

  \item Centering chirp time: To estimate a center of the peak that is
    independent of the shape of the peak, I descended both the left and right
    flanks until I reached $0.1 \times \max{\text{EOD}f_{\text{snip}}}$. The
    position of the chirp was then defined as the center of the two points on
    the flanks. The snippet was then rolled to center the chirp position within
    the snippet. As the edges could now be non-zero again, I applied another
    Tukey window ($\alpha=0.3$) to scale down fluctuations at the edges of the
    snippet. The result is a zero-centered estimate of how the frequency of the
    EOD of the fish evolves over one chirp with EOD-period resolution.

\end{enumerate}

This was repeated for each chirp for six of the recordings in the competition
dataset, in which the $\Delta\text{EOD}f$ was sufficiently large and resulted
in a total of 615 chirps. Each snippet was linearly interpolated to the
original sampling rate of \SI{20}{\kilo\hertz}.

\newpage

\begin{center}
  \textsc{Extracting chirps for hybrid simulations}
\end{center}

The following steps were taken to extract the instantaneous frequency of the
chirps for every recording in the chirp-chamber dataset and parts of the Oboti
dataset:

\begin{enumerate}

  \item Signal preprocessing: To find the EOD$f$ of the fish I used peak
    detection on the power spectrum between 300 and 1200 \si{\hertz}. The
    signal was then band-pass filtered between $\text{EOD}f - 100\si{\hertz}$
    and $\text{EOD}f + 1.9\text{EOD}f$. to remove all harmonics of the EOD.

  \item Instantaneous frequency and envelope: To estimate the instantaneous
    frequency of the EOD, I used the zero crossings of the signal to get the
    start and stop times of each period of the EOD. The instantaneous frequency
    was then computed as the inverse of the difference of the start and stop
    times of each period. The envelope of the EOD was estimated by taking the
    maximum value of each period of the EOD. Both the instantaneous frequency
    and the envelope were then linearly interpolated to the original sample
    rate of the signal.

  \item Smoothing and detection: I smoothed the instantaneous frequency to
    different degrees depending on the task: To detect peaks in the
    instantaneous frequency as chirps, I used a Gaussian kernel with a standard
    deviation of 30 samples. Peaks with a prominence between 20 and 400
    \si{\hertz} as well as a width between 4 and 4000 samples were then
    detected using the \texttt{find\_peaks} function of the \texttt{scipy}
    package \parencite{virtanenSciPyFundamentalAlgorithms2020}. I then
    extracted the snippet of the centered instantaneous frequency 200 ms before
    and after the peak. The result is a 2D matrix were each row is a snippet of
    the smoothed instantaneous frequency that contains a chirp. The same
    applies to the envelopes. Meanwhile, I always also kept the original
    instantaneous frequency and envelope.

  \item Shifting and normalization: To make all instantaneous frequencies
    relative to the EOD$f$ of the emitter, I shifted them to zero by
    subtracting the first element of each row. I then normalized the
    envelope by dividing it by the maximum value of each row. In the following,
    the detected chirps passed a series of filters to remove false positives.

  \item Filtering:
    The following filters were applied to the detected chirps:
    \begin{enumerate}

      \item Median envelope: I removed all chirps that had a median envelope
        lower than 0.2. Low median amplitudes indicate that the amplitude
        modulation in the window contains a peak instead of a trough.

      \item Frequency peak: I removed all chirps that had a peak frequency of
        more than 550 \si{\hertz}.

      \item Frequency trough: While an undershoot of the frequency does often
        happen after a chirp, it does not usually go further than 100
        \si{\hertz} below the baseline frequency. I removed all chirps that
        had a trough frequency of less than 200 \si{\hertz}.

      \item Baseline frequency: To look at how stable the baseline is in the
        chirp snippet, I looked at the first and last 100 \si{\milli\second}
        of each chirp, which should be the baseline frequency and mostly zero.
        If the 5th percentile of this part of the snippet was lower than
        \SI{-30}{\hertz}, I removed the chirp.

      \item Baseline envelope: If the same window on the normalized envelope
        had standard deviation of more than 0.2 I removed the chirp, as this
        would correspond to a very noisy signal.

      \item Envelope trough: I removed all chirps that had a trough of the
        envelope of less than 0.1 because the lower the envelope of the signal,
        the less precise the instantaneous frequency estimate. I encountered
        many cases where peaks in the instantaneous frequency where just
        artifacts that appeared when the amplitude of the signal approached
        zero.

      \item Chirp width: If the width of the chirp was less than 9
        \si{\milli\second}, I removed the chirp.

      \item If the width to height ratio was less than
        \SI{0.001}{\second\per\hertz}, I removed the chirp.

    \end{enumerate}

  \item Centering time: To find the center of the chirp, I computed a Boolean
    array that was true for all values above 20 \si{\hertz} and false for all
    values below 20 \si{\hertz}. If the number of peaks in this array was more
    than 1, or if the peak reached up to one edge of the snippet, I removed
    the chirp. I then determined the center of the chirp as the center between
    the two threshold crossings. To center the snippet, I then rolled the
    snippet so that the center of the chirp was at the center of the snippet.

\end{enumerate}
